\source{gilbarg-trudinger:page-0044}

The maximum principle is an important feature of second order elliptic equations that distinguishes them from equations of higher order and systems of equations. In addition to its many applications, the maximum principle provides pointwise estimates that lead to a more definitive theory than would be otherwise available. In this chapter, most of the results will be based solely on the ellipticity of $L$ and not on other special properties of the coefficients (such as smoothness). It is this generality which makes the maximum principle useful in apriori estimation of solutions, especially in nonlinear problems.

\subsection*{3.1. The Weak Maximum Principle}

For many purposes it suffices to have the following \emph{weak maximum principle}.

\begin{theorem}\label{gt-ch3-wmp}
\dcref{Theorem 3.1}\group{ch03-classical-maximum-principle}\source{gilbarg-trudinger:page-0044}
\emph{Let $L$ be elliptic in the bounded domain $\Omega$. Suppose that}
\begin{equation}
Lu \ge 0 \ (\le 0) \ \text{in } \Omega, \qquad c = 0 \text{ in } \Omega,
\tag{3.4}
\end{equation}
\emph{with $u \in C^2(\Omega) \cap C^0(\overline{\Omega})$. Then the maximum (minimum) of $u$ in $\overline{\Omega}$ is achieved on $\partial\Omega$, that is,}
\begin{equation}
\sup_{\Omega} u = \sup_{\partial\Omega} u \ (\inf_{\Omega} u = \inf_{\partial\Omega} u).
\tag{3.5}
\end{equation}
\end{theorem}

It is apparent that the conclusion remains valid if $|b^i|/\lambda$ is only locally bounded in $\Omega$, for example if $a^{ij},\, b^i \in C^0(\Omega)$. Also, if $u$ is not assumed continuous in $\overline{\Omega}$, the conclusion (3.5) can be replaced by

\begin{equation}
\sup_{\Omega} u = \lim_{x \to \partial\Omega} \sup u(x) \, (\inf_{\Omega} u = \lim_{x \to \partial\Omega} \inf u(x))
\tag{3.6}
\end{equation}

\begin{proof}
It is readily seen that if $Lu > 0$ in $\Omega$, then a strong maximum principle holds; that is, $u$ cannot achieve an interior maximum in $\overline{\Omega}$. For at such a point $x_0$, $Du(x_0)=0$ and the matrix $D^2u(x_0)=[D_{ij}u(x_0)]$ is nonpositive. But the matrix $[a^{ij}(x_0)]$ is positive since $L$ is elliptic. Consequently $Lu(x_0)=a^{ij}(x_0)D_{ij}u(x_0)\le 0$, contradicting $Lu>0$. (Note that only the semi-definiteness of the coefficient matrix $[a_{ij}]$ is needed in this argument.)

By hypothesis (3.3), $|b^i|/\lambda \le b_0 = \text{constant}$. Then since $a^{11}\ge \lambda$, there is a sufficiently large constant $\gamma$ for which

\[
L\,e^{\gamma x_1} = (\gamma^2 a^{11} + \gamma b^1)e^{\gamma x_1} \ge \lambda(\gamma^2-\gamma b_0)e^{\gamma x_1} > 0.
\]

Hence for any $\epsilon>0$, $L(u+\epsilon e^{\gamma x_1})>0$ in $\Omega$ so that

\[
\sup_{\Omega}(u+\epsilon e^{\gamma x_1}) = \sup_{\partial\Omega}(u+\epsilon e^{\gamma x_1})
\]

by the above. Letting $\epsilon \to 0$, we see that $\sup_{\Omega}u=\sup_{\partial\Omega}u$, as asserted in the theorem. $\square$
\end{proof}

\source{gilbarg-trudinger:page-0045}

\begin{quote}
\textit{Remark.} It is clear from the proof that the theorem holds under the weaker hypothesis that the coefficient matrix $[a^{ij}]$ is non-negative and that for some $k$ the ratio $\lvert b^k\rvert/a^{kk}$ is locally bounded.
\end{quote}

It is convenient to introduce the following terminology suggested by the maximum principle: a function satisfying $Lu=0$ $(\geq 0,\ \leq 0)$ in $\Omega$ is a \textit{solution} (\textit{subsolution}, \textit{supersolution}) of $Lu=0$ in $\Omega$. When $L$ is the Laplacian, these terms correspond respectively to harmonic, subharmonic and superharmonic functions.

Let us suppose more generally that $c\leq 0$ in $\Omega$. By considering the subset $\Omega^+ \subset \Omega$ in which $u>0$, one sees that if $Lu\geq 0$ in $\Omega$, then $L_0u=a^{ij}D_{ij}u+b^iD_iu\geq -cu\geq 0$ in $\Omega^+$ and hence the maximum of $u$ on $\overline{\Omega^+}$ must be achieved on $\partial\Omega^+$ and hence also on $\partial\Omega$. Thus, writing $u^+=\max(u,0)$, $u^-=\min(u,0)$ we obtain:

\begin{corollary}\label{gt-ch3-cor-weak}
\dcref{Corollary 3.2}\group{ch03-classical-maximum-principle}\source{gilbarg-trudinger:page-0045}
\textit{Let $L$ be elliptic in the bounded domain $\Omega$. Suppose that in $\Omega$}
\[
Lu\geq 0\ (\leq 0),\qquad c\leq 0,
\tag{3.7}
\]
\textit{with $u\in C^0(\overline{\Omega})$. Then}
\[
\sup_{\Omega} u\leq \sup_{\partial\Omega} u^+\ (\inf_{\Omega} u\geq \inf_{\partial\Omega} u^-).
\tag{3.8}
\]
\textit{If $Lu=0$ in $\Omega$, then}
\[
\sup_{\Omega} \lvert u\rvert=\sup_{\partial\Omega} \lvert u\rvert.
\tag{3.9}
\]
\end{corollary}

In this corollary, the condition $c\leq 0$ cannot be relaxed in general to allow $c>0$, as is evident from the existence of positive eigenvalues $\kappa$ for the problem:
\[
\Delta u+\kappa u=0,\quad u=0\ \text{on}\ \partial\Omega.
\]
An immediate and important application of the weak maximum principle is to the problem of uniqueness and continuous dependence of solutions on their boundary values. From Corollary 3.2 follows automatically a uniqueness result for the classical Dirichlet problem for operators $L$, and a comparison principle, which is the typical form of application of the corollary.

\begin{theorem}\label{gt-ch3-comparison}
\dcref{Theorem 3.3}\group{ch03-classical-maximum-principle}\source{gilbarg-trudinger:page-0045}
\textit{Let $L$ be elliptic in $\Omega$ with $c\leq 0$ in $\Omega$. Suppose that $u$ and $v$ are functions in $C^2(\Omega)\cap C^0(\overline{\Omega})$ satisfying $Lu = Lv$ in $\Omega$, $u = v$ on $\partial\Omega$. Then $u = v$ in $\Omega$. If $Lu\geq Lv$ in $\Omega$ and $u\leq v$ on $\partial\Omega$, then $u\leq v$ in $\Omega$.}
\end{theorem}

\section*{3.2. The Strong Maximum Principle}

Although the weak maximum principle suffices for most applications, it is often necessary to have the strong form which excludes the assumption of a non-trivial interior maximum. We shall obtain such a result for locally uniformly elliptic operators by means of the following frequently useful boundary point lemma. The domain $\Omega$ is said to satisfy an \textit{interior sphere condition} at $x_0\in \partial\Omega$ if there exists a ball $B\subset \Omega$ with $x_0\in \partial B$, (that is, the complement of $\Omega$ satisfies an \textit{exterior sphere condition} at $x_0$).

\source{gilbarg-trudinger:page-0046}

\begin{lemma}\label{gt-ch3-hopf}
\dcref{Lemma 3.4}\group{ch03-classical-maximum-principle}\source{gilbarg-trudinger:page-0046}
\emph{Suppose that $L$ is uniformly elliptic, $c=0$ and $Lu\geq 0$ in $\Omega$. Let $x_0\in\partial\Omega$ be such that}
\begin{enumerate}
\item[(i)] $u$ is continuous at $x_0$;
\item[(ii)] $u(x_0)>u(x)$ for all $x\in\Omega$;
\item[(iii)] $\partial\Omega$ satisfies an interior sphere condition at $x_0$.
\end{enumerate}
\emph{Then the outer normal derivative of $u$ at $x_0$, if it exists, satisfies the strict inequality}
\[
\tag{3.10}
\frac{\partial u}{\partial\nu}(x_0)>0.
\]
\emph{If $c\leq 0$ and $c/\lambda$ is bounded, the same conclusion holds provided $u(x_0)\geq 0$, and if $u(x_0)=0$ the same conclusion holds irrespective of the sign of $c$.}
\end{lemma}

\begin{proof}
Since $\Omega$ satisfies an interior sphere condition at $x_0$, there exists a ball $B=B_R(y)\subset\Omega$ with $x_0\in\partial B$. For $0<\rho<R$, we introduce an auxiliary function $v$ by defining
\[
v(x)=e^{-\alpha r^2}-e^{-\alpha R^2}
\]
where $r=\lvert x-y\rvert>\rho$ and $\alpha$ is a positive constant yet to be determined. Direct calculation gives for $c\leq 0$
\[
Lv(x)=e^{-\alpha r^2}\bigl[4\alpha^2 a^{ij}(x_i-y_i)(x_j-y_j)-2\alpha(a^{ii}+b^i(x_i-y_i))\bigr]+cv
\]
\[
\geq e^{-\alpha r^2}\bigl[4\alpha^2\lambda(x)r^2-2\alpha(a^{ii}+\lvert b\rvert r)+c\bigr],\qquad b=(b^1,\ldots,b^n).
\]

\emph{By assumption $a^{ij}/\lambda$, $\lvert b\rvert/\lambda$, and $c/\lambda$ are bounded. Hence $\alpha$ may be chosen large enough so that $Lv\geq 0$ throughout the annular region $A=B_R(y)-B_\rho(y)$. Since $u-u(x_0)<0$ on $\partial B_\rho(y)$ there is a constant $\epsilon>0$ for which $u-u(x_0)+\epsilon v\leq 0$ on $\partial B_\rho(y)$. This inequality is also satisfied on $\partial B_R(y)$ where $v=0$. Thus we have $L(u-u(x_0)+\epsilon v)\geq -cu(x_0)\geq 0$ in $A$, and $u-u(x_0)+\epsilon v\leq 0$ on $\partial A$. The weak maximum principle (Corollary 3.2) now implies that $u-u(x_0)+\epsilon v\leq 0$ throughout $A$. Taking the normal derivative at $x_0$, we obtain, as required,}
\[
\frac{\partial u}{\partial\nu}(x_0)\geq -\epsilon \frac{\partial v}{\partial\nu}(x_0)=-\epsilon v'(R)>0.
\]

\emph{For $c$ of arbitrary sign, if $u(x_0)=0$ the preceding argument remains valid if $L$ is replaced everywhere by $L-c^+$.} $\square$
\end{proof}

\medskip

\emph{More generally, whether or not the normal derivative exists, we get}
\[
\tag{3.11}
\liminf_{x\to x_0}\frac{u(x_0)-u(x)}{\lvert x-x_0\rvert}>0,
\]
\emph{where the angle between the vector $x_0-x$ and the normal at $x_0$ is less than $\pi/2-\delta$ for some fixed $\delta>0$.}

\source{gilbarg-trudinger:page-0047}

\section*{3.2. The Strong Maximum Principle}

Although the interior sphere condition can be relaxed somewhat, it is not possible to assert (3.11) without suitable smoothness of $\partial\Omega$ at $x_0$. For example, let $L=\Delta$ and $\Omega \subset \cR^2$ be the region in the right-half-plane in which $u=\operatorname{Re}(z/\log z)<0$. An elementary calculation shows that $\partial\Omega \subset C^1$ near $z=0$ and $u_x(0,0)=0$, so that (3.11) is false.

We are now in a position to derive the following \emph{strong maximum principle} of E. Hopf [HO 1].

\begin{theorem}\label{gt-ch3-smp}
\dcref{Theorem 3.5}\group{ch03-classical-maximum-principle}\source{gilbarg-trudinger:page-0047}
Let $L$ be uniformly elliptic, $c=0$ and $Lu\ge 0(\le 0)$ in a domain $\Omega$ (not necessarily bounded). Then if $u$ achieves its maximum (minimum) in the interior of $\Omega$, $u$ is a constant. If $c\le 0$ and $c/\lambda$ is bounded, then $u$ cannot achieve a non-negative maximum (non-positive minimum) in the interior of $\Omega$ unless it is constant.
\end{theorem}

The conclusion obviously remains valid if $L$ is only locally uniformly elliptic and $|b^i|/\lambda$, $c/\lambda$ are only locally bounded.

\begin{proof}
If we assume, contrary to the theorem, that $u$ is non-constant and achieves its maximum $M\ge 0$ in the interior of $\Omega$, then the set $\Omega^-$ on which $u<M$ satisfies $\Omega^- \subset \Omega$ and $\partial\Omega^- \cap \Omega \ne \phi$. Let $x_0$ be a point in $\Omega^-$ that is closer to $\partial\Omega^-$ than to $\partial\Omega$, and consider the largest ball $B\subset \Omega^-$ having $x_0$ as center. Then $u(y)=M$ for some point $y\in \partial B$ while $u<M$ in $B$. The preceding lemma implies that $Du(y)\ne 0$, which is impossible at the interior maximum $y$. $\square$
\end{proof}

If $c<0$ at some point, then the constant of the theorem is obviously zero. Also, if $u=0$ at an interior maximum (minimum), then it follows from the proof of the theorem that $u\equiv 0$, irrespective of the sign of $c$.

It is of course possible to prove the strong maximum principle directly without going through Theorem 3.1 and Lemma 3.4; (see [MR 2] for example).

Uniqueness theorems for other types of boundary value problems are consequences of Lemma 3.4 and Theorem 3.5. In particular, we have the following uniqueness theorem for the classical Neumann problem.

\begin{theorem}\label{gt-ch3-neumann}
\dcref{Theorem 3.6}\group{ch03-classical-maximum-principle}\source{gilbarg-trudinger:page-0047}
\emph{Let $u\in C^2(\Omega)\cap C^0(\bar{\Omega})$ be a solution of $Lu=0$ in the bounded domain $\Omega$, where $L$ is uniformly elliptic, $c\le 0$, $c/\lambda$ is bounded and $\Omega$ satisfies an interior sphere condition at each point of $\partial\Omega$. If the normal derivative is defined everywhere on $\partial\Omega$ and $\partial u/\partial\nu=0$ on $\partial\Omega$, then $u$ is constant in $\Omega$. If, also, $c<0$ at some point in $\Omega$, then $u\equiv 0$.}
\end{theorem}

\begin{proof}
If $u\ne \mathrm{const.}$, we may assume that either of the functions $u$ or $-u$ achieves a non-negative maximum $M$ at a point $x_0$ on $\partial\Omega$ and is less than $M$ in $\Omega$ (by the strong maximum principle). Applying Lemma 3.4 at $x_0$ we infer that $\partial u/\partial\nu(x_0)\ne 0$, contradicting the hypothesis. $\square$
\end{proof}

The result of Theorem 3.6 may also be generalized to mixed boundary value and oblique derivative problems; (see Problem 3.1). When $\partial\Omega$ has corners or edges where the derivatives of $u$ are not defined, these results are false in the stated generality, even if $u$ is assumed continuous on $\bar{\Omega}$; (see Problem 3.8(a)).

\source{gilbarg-trudinger:page-0048}

\section*{3.3. Apriori Bounds}

The maximum principle also provides a simple pointwise estimate for solutions of the inhomogeneous equation $Lu=f$ in bounded domains. We remark that only the ellipticity and bounds on the coefficients are involved. This proves to be an important consideration in nonlinear problems.

\begin{theorem}\label{gt-ch3-apriori}
\dcref{Theorem 3.7}\group{ch03-classical-maximum-principle}\source{gilbarg-trudinger:page-0048}
\textit{Let $Lu\ge f$ ($=f$) in a bounded domain $\Omega$, where $L$ is elliptic, $c\le 0$, and $u\in C^0(\overline{\Omega})\cap C^2(\Omega)$. Then}
\begin{equation}
\tag{3.12}
\sup_{\Omega} u^-\le \sup_{\partial\Omega} u^+ + C \sup_{\Omega}\frac{|f^-|}{\lambda},
\end{equation}
\noindent\textit{where $C$ is a constant depending only on diam $\Omega$ and $\beta=\sup |b|/\lambda$. In particular, if $\Omega$ lies between two parallel planes a distance $d$ apart, then (3.12) is satisfied with $C=e^{(\beta+1)d}-1$.}
\end{theorem}

\begin{proof}
Let $\Omega$ lie in the slab $0<x_1<d$, and set $L_0=a^{ij}D_{ij}+b^iD_i$. For $\alpha\ge \beta+1$ we have
\[
L_0e^{\alpha x_1}=(\alpha^2a^{11}+\alpha b^1)e^{\alpha x_1}\ge \lambda(\alpha^2-\alpha\beta)e^{\alpha x_1}\ge \lambda.
\]
Let
\[
v=\sup_{\partial\Omega}u^+ + (e^{\alpha d}-e^{\alpha x_1})\sup_{\Omega}\frac{|f^-|}{\lambda}.
\]
Then, since $Lv=L_0v+cv\le -\lambda\sup_{\partial\Omega}(|f^-|/\lambda)$,
\[
L(v-u)\le -\lambda\left(\sup_{\Omega}\frac{|f^-|}{\lambda}+\frac{f}{\lambda}\right)\le 0 \quad \text{in }\Omega,
\]
and $v-u\ge 0$ on $\partial\Omega$. Hence, for $C=e^{\alpha d}-1$ and $\alpha\ge \beta+1$, we obtain the desired result for the case $Lu\ge f$,
\[
\sup_{\Omega}u\le \sup_{\Omega}v\le \sup_{\partial\Omega}u^+ + C\sup_{\Omega}\frac{|f^-|}{\lambda}.
\]
Replacing $u$ by $-u$, we obtain (3.12) for the case $Lu=f$. $\square$
\end{proof}

Theorem 3.7 will be strengthened in Chapters 8 and 9 to provide analogous estimates for $\sup u$ in terms of integral norms of $f$.

\noindent \textbf{When the condition $c\le 0$ is not satisfied, it is still possible to assert an apriori bound analogous to (3.12) provided the domain $\Omega$ lies between sufficiently close parallel planes.}

\source{gilbarg-trudinger:page-0049}

\section*{3.4. Gradient Estimates for Poisson's Equation}

\begin{corollary}\label{gt-ch3-grad-cor}
\dcref{Corollary 3.8}\group{ch03-classical-maximum-principle}\source{gilbarg-trudinger:page-0049}
\textit{Let $Lu=f$ in a bounded domain $\Omega$, where $L$ is elliptic and $u\in C^0(\overline{\Omega})\cap C^2(\Omega)$. Let $C$ be the constant of Theorem 3.7 and suppose that}
\begin{equation}
\tag{3.13}
C_1 = 1 - C \sup_{\Omega}\frac{c^+}{\lambda} > 0.
\end{equation}
\textit{Then}
\begin{equation}
\tag{3.14}
\sup_{\Omega}|u| \le \frac{1}{C_1}\left(\sup_{\partial\Omega}|u| + C \sup_{\Omega}\frac{|f|}{\lambda}\right).
\end{equation}
\end{corollary}

\noindent \textit{Remark.} Since $C=e^{(\beta+1)d}-1$ is a possible value of the constant in (3.12), where $d$ is the width of any slab containing $\Omega$, condition (3.13) will be satisfied in any sufficiently narrow domain in which the quantities $|b|/\lambda$ and $c/\lambda$ are bounded above. If $c^+\equiv 0$ (i.e., $c\le 0$), then $C_1=1$ and (3.14) reduces to (3.12).

\begin{proof}
of Corollary 3.8. Let us rewrite $Lu=(L_0+c)u=f$ in the form
\[
(L_0+c^-)u=f' \equiv f+(c^- - c)u = f - c^+u.
\]
From (3.12) we obtain
\[
\sup_{\Omega}|u| \le \sup_{\partial\Omega}|u| + C \sup_{\Omega}\frac{|f'|}{\lambda}
\]
\[
\le \sup_{\partial\Omega}|u| + C\left(\sup_{\Omega}\frac{|f|}{\lambda} + \sup_{\Omega}|u| \sup_{\Omega}\frac{c^+}{\lambda}\right).
\]
This inequality and (3.13) imply (3.14). \hfill $\square$
\end{proof}

\medskip

\noindent An immediate consequence of Corollary 3.8 is uniqueness for solutions of the Dirichlet problem in sufficiently small domains (assuming of course fixed upper bounds on the quantities $|b|/\lambda$ and $c/\lambda$).

\section*{3.4. Gradient Estimates for Poisson's Equation}

The maximum principle can also be used to derive estimates on derivatives of solutions provided additional conditions are placed on the equation. To illustrate the method we obtain such estimates for Poisson's equation. The results derived here will not be required for later developments.

\medskip

\noindent Let $\Delta u=f$ in the cube $Q=\{x=(x_1,\dots,x_n)\in \mathbb{R}^n \mid |x_i|<d,\, i=1,\dots,n\}$, with $u\in C^2(Q)\cap C^0(\overline{Q})$ and $f$ bounded in $Q$. By means of a comparison argument we shall derive the estimate
\begin{equation}
\tag{3.15}
|D_i u(0)| \le \frac{n}{d}\sup_{\partial Q}|u| + \frac{d}{2}\sup_Q|f|,\qquad i=1,\dots,n.
\end{equation}

\source{gilbarg-trudinger:page-0050}

\section*{3. The Classical Maximum Principle}

In the half-cube
\[
Q'=\{(x_1,\dots,x_n)\mid |x_i|<d,\ i=1,\dots,n-1,\ 0<x_n<d\},
\]
consider the function
\[
\varphi(x',x_n)=\frac{1}{2}[u(x',x_n)-u(x',-x_n)],
\]
where we write $x'=(x_1,\dots,x_{n-1})$ and $x=(x',x_n)$. One sees that $\varphi(x',0)=0$,
\[
\sup_{\partial Q'}|\varphi|\le M=\sup_{\partial Q}|u|,\qquad \text{and}\qquad |\Delta\varphi|\le N=\sup_Q|f|
\]
in $Q'$. Consider also the function
\[
\psi(x',x_n)=\frac{M}{d^2}\bigl[|x'|^2+x_n(nd-(n-1)x_n)\bigr]+N\frac{x_n}{2}(d-x_n).
\]
Obviously $\psi(x',x_n)\ge 0$ on $x_n=0$ and $\psi\ge M$ on the remaining portion of $\partial Q'$; also $\Delta\psi=-N$. Hence $\Delta(\psi\pm\varphi)\le 0$ in $Q'$ and $\psi\pm\varphi\ge 0$ on $\partial Q'$, from which it follows by the maximum principle that $|\varphi(x',x_n)|\le\psi(x',x_n)$ in $Q'$. Letting $x'=0$ in the expressions for $\varphi$ and $\psi$, then dividing by $x_n$ and letting $x_n$ tend to zero, we obtain
\[
|D_nu(0)|=\lim_{x_n\to 0}\left|\frac{\varphi(0,x_n)}{x_n}\right|\le \frac{n}{d}M+\frac{d}{2}N,
\]
which is the asserted estimate (3.15) for $i=n$. The result follows in the corresponding way for $i=1,\dots,n-1$. If $f=0$, (3.15) provides an independent proof of (essentially) the gradient bound (2.31) for harmonic functions. From (3.15) we infer that in any domain $\Omega$ a bounded solution $u$ of $\Delta u=f$ satisfies an estimate
\[
(3.16)\qquad \sup_{\Omega} d_x|Du(x)|\le C\left(\sup_{\Omega}|u|+\sup_{\Omega} d_x^2|f(x)|\right),
\]
where $d_x=\operatorname{dist}(x,\partial\Omega)$ and $C=C(n)$. For if $x\in\Omega$ and $Q$ is a cube of side $d=d_x/\sqrt{n}$ with its center at $x$, we have from (3.15) the inequality
\[
d_x|Du(x)|\le C\left(\sup_{\partial Q}|u|+d^2\sup_Q|f|\right)
\]
\[
\le C\left(\sup_{\Omega}|u|+\sup_{\Omega} d_y^2|f(y)|\right).
\]
(Here we have used the same letter $C$ to denote constants depending only on $n$.) In the same generality as above, we now derive by a similar comparison argument an estimate of the modulus of continuity of the gradient of solutions of Poisson's equation.

\source{gilbarg-trudinger:page-0051}

\begin{quote}
\textbf{3.4. Gradient Estimates for Poisson's Equation}
\end{quote}

Again let $u \in C^2(Q) \cap C^0(\overline{Q})$ be a solution of $\Delta u = f$ in the cube $Q$, and set $M=\sup_Q |u|$, $N=\sup_Q |f|$. Let $Q'$ be the domain in $\mathbb{R}^{n+1}$ given by
\[
Q'=\{(x_1,\ldots,x_{n-1},y,z)\mid |x_i|<d/2,\, i=1,\ldots,n-1,\, 0<y,z<d/4\},
\]
and let us define in $Q'$ the function
\[
\phi(x',y,z)=\frac14\bigl[u(x',y+z)-u(x',y-z)-u(x',-y+z)+u(x',-y-z)\bigr].
\]
Introducing the elliptic operator
\[
L=\sum_{i=1}^{n-1}\frac{\partial^2}{\partial x_i^2}+\frac12\frac{\partial^2}{\partial y^2}+\frac12\frac{\partial^2}{\partial z^2}
\]
in the $n+1$ variables $x_1,\ldots,x_{n-1},y,z$, we see that $|L\phi|\leq N$ in $Q'$. Also, on $Q'$ we have: (i) $\phi(x',0,z)=\phi(x',y,0)=0$; (ii) $|\phi|\leq M$ on $|x_i|=d/2$, $i=1,\ldots,n-1$; (iii) $|\phi(x',d/4,z)|\leq \mu z$ and $|\phi(x',y,d/4)|\leq \mu y$, where $|Du|\leq \mu$ in $Q'$, $\mu$ being given in terms of $M$ and $N$ by (3.16). We now choose a comparison function in $Q'$ of the form,
\[
\tag{3.17}
\psi(x',y,z)=\frac{4M|x'|^2}{d^2}+\frac{4\mu}{d}yz+kyz\log\frac{2d}{y+z},
\]
where $k$ is a positive constant yet to be determined. We observe first that $|\phi|\leq \psi$ on $\partial Q'$. Since
\[
L\psi=\frac{8(n-1)}{d^2}M+k\left(-1+\frac{yz}{(y+z)^2}\right)\leq \frac{8(n-1)M}{d^2}-\frac34k,
\]
we see that $L\psi\leq -N$ provided
\[
k\geq \frac43\left(N+\frac{8(n-1)M}{d^2}\right).
\]
With such a choice of $k$, the function
\[
\psi(x',y,z)=\frac{4M|x'|^2}{d^2}+yz\left(\frac{4\mu}{d}+k\log\frac{2d}{y+z}\right)
\]
satisfies the conditions, $L(\psi\pm\phi)\leq 0$ in $Q'$, $\psi\pm\phi\geq 0$ on $\partial Q'$. Accordingly, $|\phi|\leq \psi$ in $Q'$. Letting $x'=0$ in this inequality, then dividing by $z$ and letting $z$ tend to zero, we obtain
\[
\tag{3.18}
\frac12|u_y(0,y)-u_y(0,-y)|=\lim_{z\to 0}\frac{|\phi(0,y,z)|}{z}\leq \frac{4\mu}{d}y+ky\log\frac{2d}{y}.
\]

\source{gilbarg-trudinger:page-0052}

With a slight modification of the argument an analogous estimate can be derived for $|D_i u(0,x_n)-D_i u(0,-x_n)|$ (where $D_i=\partial/\partial x_i$, $i=1,\ldots,n-1$). Let us define
\[
\varphi(\hat{x},y,z)=\frac{1}{4}\bigl[u(\hat{x},y,z)-u(\hat{x},-y,z)-u(\hat{x},y,-z)+u(\hat{x},-y,-z)\bigr]
\]
where $\hat{x}=(x_1,\ldots,x_{n-2})$. In the domain
\[
Q'=\{(x_1,\ldots,x_{n-2},y,z)\mid |x_i|<d/2,i=1,\ldots,n-2,0<y,z<d/2\}
\]
we choose a comparison function similar to (3.17) of the form
\[
\psi(\hat{x},y,z)=\frac{4M|\hat{x}|^2}{d^2}+yz\left(\frac{4\mu}{d}+\bar{k}\log\frac{2d}{y+z}\right),
\]
where $\mu$ and $\bar{k}$ are constants such that $|Du|\le \mu$ in $Q'$ and
\[
\bar{k}\ge \frac{2}{3}[N+8(n-2)M/d^2].
\]
One verifies that $\Delta(\psi\pm\varphi)\le 0$ in $Q'$ and $\psi\pm\varphi\ge 0$ on $\partial Q'$, from which it follows $|\varphi|\le \psi$ in $Q'$. As above, if we set $\hat{x}=0$ in this inequality, then divide by $y$ and let $y$ tend to zero, we obtain
\begin{equation}
\tag{3.19}
\frac{1}{2}|D_{n-1}u(0,z)-D_{n-1}u(0,-z)|=\lim_{y\to 0}\frac{|\varphi(0,y,z)|}{y}\le \frac{4\mu}{d}z+\bar{k}z\log\frac{2d}{z}.
\end{equation}
Obviously the same result is obtained if $D_{n-1}$ is replaced by $D_i$, $i=1,\ldots,n-2$.
We note that unlike the argument for (3.18), the proof of (3.19) did not require the introduction of an operator in $\R^{n+1}$.

If now $\Delta u=f$ in a domain $\Omega$ of $\R^n$, we can obtain from (3.18) and (3.19) an estimate for $|Du(x)-Du(y)|$, where $x$ and $y$ are any two points of $\Omega$. Let $d_x=\operatorname{dist}(x,\partial\Omega)$, $d_y=\operatorname{dist}(y,\partial\Omega)$ and $d_{x,y}=\min(d_x,d_y)$. Assume $d_x\le d_y$, so that $d_x=d_{x,y}$. Suppose first that $|x-y|\le d=d_x/(2\sqrt{n})$, and consider the segment joining $x$ and $y$. We choose the center of this segment as origin and rotate coordinates so that $x$ and $y$ lie on the $x_n$ axis with $x=(0,x_n)$, $y=(0,-x_n)$ in the new coordinates. The cube $Q=\{(x_1,\ldots,x_n)\mid |x_i|<d,\, i=1,\ldots,n\}$ lies in $\Omega$ at a distance greater than $d_x/2$ from $\partial\Omega$. We may apply (3.16), (3.18) and (3.19) directly in $Q$ to obtain
\[
d^2\frac{|Du(x)-Du(y)|}{|x-y|}\le C\left(\sup_Q |u|+d^2\sup_Q |f|\right)\log\frac{2d}{|x-y|},
\]
for some constant $C=C(n)$. Hence
\[
d_{x,y}^2\frac{|Du(x)-Du(y)|}{|x-y|}\le C\left(\sup_{\Omega}|u|+\sup_{\Omega}d_x^2|f(x)|\right)\log\frac{d_{x,y}}{|x-y|}.
\]

\source{gilbarg-trudinger:page-0053}

If $x$ and $y$ are now points in $\Omega$ such that $\lvert x-y\rvert > d$, we have from (3.16),
\[
d_{x,y}^{2}\frac{\lvert Du(x)-Du(y)\rvert}{\lvert x-y\rvert}
\le C\bigl(\sup_{\Omega}\lvert u\rvert+\sup_{\Omega} d_{x}^{2}\lvert f(x)\rvert\bigr).
\]
Combining these results, we obtain
\[
d_{x,y}^{2}\frac{\lvert Du(x)-Du(y)\rvert}{\lvert x-y\rvert}
\le C\bigl(\sup_{\Omega}\lvert u\rvert+\sup_{\Omega} d_{x}^{2}\lvert f(x)\rvert\bigr)
\left(\log\frac{d_{x,y}}{\lvert x-y\rvert}+1\right)
\]
where $C$ is a constant depending only on $n$.
The preceding results are collected in the following.

\begin{theorem}\label{gt-ch3-poisson}
\dcref{Theorem 3.9}\group{ch03-classical-maximum-principle}\source{gilbarg-trudinger:page-0053}
\textit{Let $u \in C^2(\Omega)$ satisfy Poisson's equation, $\Delta u = f$, in $\Omega$. Then}
\[
\sup_{\Omega} d_{x}\lvert Du(x)\rvert \le C\bigl(\sup_{\Omega}\lvert u\rvert+\sup_{\Omega} d_{x}^{2}\lvert f(x)\rvert\bigr),
\]
\textit{and for all $x, y$ in $\Omega$, $x \neq y$,}
\[
(3.20)\qquad
d_{x,y}^{2}\frac{\lvert Du(x)-Du(y)\rvert}{\lvert x-y\rvert}
\le C\bigl(\sup_{\Omega}\lvert u\rvert+\sup_{\Omega} d_{x}^{2}\lvert f(x)\rvert\bigr)
\left(\log\frac{d_{x,y}}{\lvert x-y\rvert}+1\right),
\]
\textit{where $C = C(n)$. Here $d_x = \operatorname{dist}(x,\partial\Omega)$, $d_{x,y} = \min(d_x,d_y)$.}
\end{theorem}

Despite the elementary character of its proof, this theorem is essentially sharp and the estimate (3.20) cannot be improved without further continuity assumptions on $f$. Theorem 3.9 will also hold for weak solutions in the sense of Chapter 8 provided $f$ is bounded; (see Problem 8.4). Extensions of the above results for the case of H\"older continuous $f$ are treated by other methods in Chapter 4, although the comparison methods of this section can be used to obtain these extensions as well; (see [BR 1, 2]).

\section*{3.5. A Harnack Inequality}

The maximum principle provides an elementary proof of a general Harnack inequality for uniformly elliptic equations in two independent variables. Letting $D_{\rho}=D_{\rho}(0)$ denote the open disk of radius $\rho$ centered at the origin, we state the result in the following form.

\begin{theorem}\label{gt-ch3-harnack}
\dcref{Theorem 3.10}\group{ch03-classical-maximum-principle}\source{gilbarg-trudinger:page-0053, gilbarg-trudinger:page-0054, gilbarg-trudinger:page-0055, gilbarg-trudinger:page-0056}
Let $u$ be a non-negative $C^2$ solution of
\[
Lu=au_{xx}+2bu_{xy}+cu_{yy}=0
\]
\source{gilbarg-trudinger:page-0054}

in the disk $D_R$, and suppose that $L$ is uniformly elliptic in $D_R$. Then at all points $z=(x, y)\in D_{R/4}$ we have the inequality
\[
(3.21)\qquad Ku(0)\le u(z)\le K^{-1}u(0),
\]
\textit{where $K$ is a constant depending only on the ellipticity modulus $\mu=\sup_D \Lambda/\lambda$.}
\end{theorem}

\begin{proof}
We note to begin with that since the equation $Lu=0$ and the modulus $\mu$ are invariant under similarity transformation, it suffices to prove the theorem in the unit disk $D=D_1$. Since $u\ge 0$ in $D$, the strong maximum principle (Theorem 3.5) implies that either $u\equiv 0$ or $u>0$ in $D$, so it suffices to assume the latter. Consider the set $G$ in $D$ where $u>u(0)/2$, and let $G'\subset G$ be the component containing $0$. One sees from the maximum principle that $\partial G'\cap \partial D$ is non-empty, and hence there is no loss of generality in assuming that the point $Q=(0,1)$ is in $\partial G'$. We define the functions $v_+$ and $v_-$ by
\[
v_{\pm}(x, y)=\pm x+\frac{3}{4}-k\left(y-\frac{1}{2}\right)^2,
\]
where $k$ is a positive constant. The parabolas, $\Gamma_{\pm} : v_{\pm}=0$, have vertices $\left(\mp \frac{3}{4}, \frac{1}{2}\right)$ in $D$ and common axis $y=\frac{1}{2}$. If $k$ is sufficiently large (it suffices that $k\ge 3$), the domains $P_{\pm}$ in $D$ in which $v_{\pm}>0$ have an intersection $P_+\cap P_-$ bounded by arcs of $\Gamma_+,\Gamma_-$ and lying entirely in the upper half of $D$; (see Figure 1). In $P_{\pm}$, the functions $v_{\pm}$ obviously satisfy the inequality $0<v_{\pm}<\frac{3}{4}$.

Figure 1

Setting $E_{\pm}=\exp(\alpha v_{\pm})$, where $\alpha$ is a positive constant yet to be chosen, we find by direct calculation
\[
LE_{\pm}=E_{\pm}\left\{\alpha^2\left[a\mp 4bk\left(y-\frac{1}{2}\right)+4ck^2\left(y-\frac{1}{2}\right)^2\right]-2\alpha kc\right\}
\]
\[
\ge E_{\pm}\left(\alpha^2\lambda-2\alpha k\Lambda\right)
\]
\[
\ge 0\text{ in }D,\quad \text{if }\alpha\ge 2k\mu.
\]

\source{gilbarg-trudinger:page-0055}

Consequently, with such a choice of $\alpha$, the functions
\[
(3.22)\qquad w_{\pm}=(E_{\pm}-1)/(e^{7\alpha/4}-1)
\]
have the properties:
\[
Lw_{\pm}\ge 0\text{ in }D;\qquad w_{\pm}=0\text{ on }\Gamma_{\pm};\qquad 0<w_{\pm}<1\text{ in }P_{\pm}.
\]

Now let $z$ be any point in $P_{+}\cap P_{-}$. Then either: (i) $u\ge u(0)/2$ and $z\in \overline{G}$; or (ii) $z$ lies in a component $U_{+}$ of $P_{+}-\overline{G}$ such that $\partial U_{+}\subset \Gamma_{+}\cup \partial G$; or (iii) $z$ lies in a component $U_{-}$ of $P_{-}-\overline{G}$ such that $\partial U_{-}\subset \Gamma_{-}\cup \partial G$; (see Figure 1). These are the only alternatives, since either $P_{+}\cap P_{-}\subset G'$ or $\partial G'$ separates $P_{+}\cup P_{-}$. (The two dimensionality is used here in an essential way.) In cases (ii) and (iii) we have
\[
u-\frac{1}{2}u(0)w_{\pm}=\frac{1}{2}u(0)(1-w_{\pm})>0 \qquad \text{on }\partial G\cap \partial U_{\pm},
\]
\[
u-\frac{1}{2}u(0)w_{\pm}=u>0 \qquad \text{on }\Gamma_{\pm}\cap \partial U_{\pm}.
\]
Thus $u-\frac{1}{2}u(0)w_{\pm}>0$ on $\partial U_{\pm}$. Since $L(u-\frac{1}{2}u(0)w_{\pm})\le 0$, we infer that
\[
u(z)>\frac{1}{2}u(0)\min\bigl(w_{+}(z),w_{-}(z)\bigr)\qquad \forall z\in P_{+}\cap P_{-}.
\]
In particular, on the segment, $y=\frac{1}{2}, |x|\le \frac{1}{2}$, we have
\[
(3.23)\qquad u(x,\tfrac{1}{2})>K_{1}u(0)\qquad \forall x\in \bigl[-\tfrac{1}{2},\tfrac{1}{2}\bigr],
\]
where
\[
K_{1}=\frac{1}{2}(e^{\alpha/4}-1)/(e^{7\alpha/4}-1)=\frac{1}{2}\inf_{|x|\le 1/2}\bigl[w_{+}(x,\tfrac{1}{2}),w_{-}(x,\tfrac{1}{2})\bigr].
\]
We now define another comparison function, similar to (3.22). Namely, setting
\[
v=y+1-6x^{2},
\]
we consider the domain
\[
P=\{(x,y)\in D\mid v(x,y)>0,\ y<\tfrac{1}{2}\}.
\]
P is bounded by the segment, $y=\frac{1}{2}, |x|\le \frac{1}{2}$, and the arc $\Gamma$ of the parabola $v=0$, with vertex at $(0,-1)$ and passing through the points $(\pm \tfrac{1}{2},\tfrac{1}{2})$. As before, for a suitable choice of $\beta>0$ depending only on $\mu$, the function
\[
w=(e^{\beta v}-1)/(e^{3\beta/2}-1)
\]
has the properties:
\[
Lw\ge 0\text{ in }D;\qquad w=0\text{ on }\Gamma;\qquad 0<w<1\text{ in }P.
\]
From (3.23) we have that $u-K_{1}u(0)w>0$ on $\partial P$, and since $L(u-K_{1}u(0)w)\le 0$, it follows from the maximum principle that
\[
u(z)>K_{1}u(0)w(z)\qquad \forall z\in P.
\]

\source{gilbarg-trudinger:page-0056}

Noting that $D_{1/3}\subset P$, and setting $K_2=\inf_{D_{1/3}} w$, we obtain
\[
(3.24)\qquad u(z)>K_1K_2u(0)=Ku(0)\qquad \forall z\in D_{1/3}.
\]
Clearly $K$ depends only on $\mu$.

If now $z\in D_{1/4}$, the disk $D_{3/4}(z)$ is contained in $D$ and the inequality (3.24), applied in the disk $D_{1/4}(z)$, implies
\[
u(0)>Ku(z)\qquad \forall z\in D_{1/4}.
\]
Combining this inequality with (3.24), we obtain
\[
Ku(0)<u(z)<K^{-1}u(0)\qquad \forall z\in D_{1/4}. \quad \square
\]
\end{proof}
It follows immediately from (3.21) that
\[
(3.25)\qquad \sup_{D_{R/4}}u\le \kappa\inf_{D_{R/4}}u,
\]
where $\kappa=1/K^2$. By the same chaining argument as in Theorem 2.5, we obtain the following Harnack inequality for arbitrary domains in $\mathbb{R}^2$.

\begin{corollary}\label{gt-ch3-harnack-cor}
\dcref{Corollary 3.11}\group{ch03-classical-maximum-principle}\source{gilbarg-trudinger:page-0056}
\textit{Let the hypotheses of Theorem 3.10 hold in a domain $\Omega\subset \mathbb{R}^2$. Then for any bounded subdomain $\Omega'\subset \subset \Omega$, there is a constant $\kappa$ depending only on $\Omega,\Omega'$ and $\mu$ such that}
\[
(3.26)\qquad \sup_{\Omega'}u\le \kappa\inf_{\Omega'}u.
\]
\end{corollary}

If we consider the more general elliptic equation
\[
(3.27)\qquad Lu=a^{ij}D_{ij}u+b^iD_i u+cu=0,\qquad c\le 0,\qquad i,j=1,2,
\]
where the coefficients of the operator $L$ are bounded and $\lambda\ge \lambda_0>0$, the proof of Theorem 3.10 in the unit disk $D$ is still valid (with slight modification); and the conclusion remains the same, but the constant $K$ now depends on the bounds for the coefficients in $D$ as well as on $\mu$. In stating the analogous result for a disk of radius $R$, the constant $K$ will therefore depend on $R$ in addition to the other quantities; (see Problem 3.4).

The Harnack inequality (3.21) has as consequence the following Liouville theorem.

\begin{corollary}\label{gt-ch3-liouville}
\dcref{Corollary 3.12}\group{ch03-classical-maximum-principle}\source{gilbarg-trudinger:page-0056, gilbarg-trudinger:page-0057}
\textit{If the equation $Lu=au_{xx}+2bu_{xy}+cu_{yy}=0$ is uniformly elliptic in $\mathbb{R}^2$ and $u$ is a solution bounded below (or above) and defined over the entire plane, then $u$ is constant.}
\end{corollary}

\source{gilbarg-trudinger:page-0057}

\section*{3.6. Operators in Divergence Form}

\begin{proof}
We may assume that inf $u=0$, and, hence, for any $\varepsilon>0$, $u(z_0)<\varepsilon$ for some $z_0$. In every disc $D_{2R}(z_0)$, we have from (3.21) that $u(z)<K\varepsilon$ for all $z\in D_R(z_0)$. Since $K$ is a constant independent of $R$, it follows that $u(z)<K\varepsilon$ for all $z\in \mathbb{R}^2$, and the conclusion is immediate by letting $\varepsilon\to 0$. \qedhere
\end{proof}

A proof of the extension of the Harnack inequality (Theorem 3.10) and of Corollary 3.12 to higher dimensions appears in Chapter 9. Other Harnack inequalities, for equations in divergence form, together with some important applications, are contained in Chapters 8 and 13.

\section*{3.6. Operators in Divergence Form}

We conclude this chapter with a brief look at the situation for operators in divergence form. In many situations it is more natural to consider these than operators of the form (3.1). The simplest such case is
\[
\tag{3.28}
Lu = D_j(a^{ij}D_i u).
\]
Later it will be necessary to consider more general operators whose principal part is in divergence form. $L$ will be called elliptic in $\Omega$ if the coefficient matrix $[a^{ij}(x)]$ is positive for all $x\in\Omega$.

Evidently the results concerning the maximum principle apply equally well to the operator (3.28) when the coefficients $a^{ij}$ are sufficiently smooth. However, when this is not the case, or, as in nonlinear problems, when it is often inappropriate to make quantitative assumptions concerning the smoothness of the coefficients (e.g., bounds on their derivatives), the essentially algebraic methods of the earlier part of this chapter cease to be applicable and must be replaced by integral methods, which are more natural for the divergence structure of $L$.

The relations $Lu=0$ ($\ge 0$, $\le 0$) satisfied by solutions (subsolutions, supersolutions) of $Lu=0$ can be defined for broader classes of coefficients and functions $u$ than those formally permitted in (3.28). Thus, if the coefficients $a^{ij}$ are bounded and measurable and $u\in C^1(\Omega)$, then, in a generalized sense, $u$ is said to satisfy $Lu=0$ ($\ge 0$, $\le 0$) respectively in $\Omega$, according as
\[
\tag{3.29}
\int_{\Omega} a^{ij}(x)D_i u\,D_j \varphi\,dx = 0\quad (\le 0,\ge 0)
\]
for all non-negative functions $\varphi\in C_0^1(\Omega)$. By application of the divergence theorem this is easily seen to be equivalent to $Lu=0$ ($\ge 0$, $\le 0$) when $a^{ij}\in C^1(\Omega)$ and $u\in C^2(\Omega)$. In later chapters generalized solutions will be defined in wider and more appropriate function spaces.

\source{gilbarg-trudinger:page-0058}

The weak maximum principle is an immediate consequence of (3.29). For let $u$ satisfy
\[
\int_{\Omega} a^{ij} D_i u\, D_j \varphi \, dx \leq 0 \qquad \text{for all } \varphi \in C_0^1(\Omega), \ \varphi \geq 0;
\tag{3.30}
\]
and suppose, contrary to our assertion, that $\sup_{\Omega} u > \sup_{\partial\Omega} u = u_0$. Then for some constant $c>0$, there is a subdomain $\Omega' \subset \subset \Omega$ in which $v = u-u_0-c > 0$ and $v = 0$ on $\partial\Omega'$. The relation (3.30) remains true with $u$ replaced by $v$ and with $\varphi = v$ in $\Omega'$, $= 0$ elsewhere. (As thus defined $\varphi \notin C_0^1(\Omega)$, but (3.30) can be seen to hold by approximating this $\varphi$ with functions in $C_0^1(\Omega)$.) It follows that
\[
\int_{\Omega'} a^{ij} D_i v\, D_j v\, dx \leq 0,
\]
and hence since $[a^{ij}]$ is positive, we infer that $D v = 0$ in $\Omega'$. Since $v = 0$ on $\partial\Omega'$, we have $v = 0$ in $\Omega'$, which contradicts the definition of $v$. This establishes the weak maximum principle.

Stronger and more general maximum principles for divergence structure operators will be presented in later chapters. Aside from the already noted difference in methods in treating the two classes of operators, it should be remarked also that results relating to the maximum principle are often different for the operators (3.1) and (3.28) when there are weak smoothness conditions on the coefficients. For example, Lemma 3.4 is not necessarily true for the uniformly elliptic operator of divergence form (3.28) even when the coefficients are arbitrarily smooth in the interior and continuous up to the boundary; (see Problem 3.9).

Notes

The boundary point lemma (Lemma 3.4) as proved here is due to Hopf [HO5]; an independent proof, differing only in the choice of comparison function, was obtained by Oleinik [OL]. The result remains valid, under the same hypotheses on the coefficients, if $\partial\Omega$ has a Dini continuous normal [KH]. A further extension, valid for a class of domains including Lipschitz domains, provides a proof of uniqueness for the Neumann problem in such domains [ND]. Lemma 3.4 is false in general for strictly and uniformly elliptic equations of divergence form even if the coefficients are continuous at the boundary point (see Problem 3.9), but is true if the coefficients are Hölder continuous in a neighborhood (Finn-Gilbarg [FG 1]). Results analogous to Lemma 3.4 for domains satisfying an interior cone condition, in place of the interior sphere condition, have been obtained by Oddson [OD] and Miller [ML 1, 3]. They prove (3.11) and more precise results, with $|x - x_0|^\mu$ in place of $|x - x_0|$, the exponent $\mu$ depending only on the cone angle and the ellipticity constant; (here the vector $x - x_0$ lies within a fixed subcone of the assumed interior cone at $x_0$). These essentially sharp results are based on the extremal elliptic operators of Pucci [PU 2].

\source{gilbarg-trudinger:page-0059}

The maximum principle in the generality of Theorem 3.5 was first proved by Hopf [HO 1]. For earlier results, under more restrictive hypotheses, see references in [PW], p. 156, where there is also a discussion of various extensions of the maximum principle. Some of these are considered in Chapters 8 and 9.

Section 3.4 is based on the ideas of Brandt [BR 1, 2], who has shown that much of the linear theory of classical solutions of second order elliptic and parabolic equations, including the deeper estimates of Chapters 4 and 6, can be derived from comparison arguments using the maximum principle. As in Section 3.4, the method requires appropriate (and generally not obvious) choices of comparison functions, which are used to estimate difference quotients and hence derivatives.

The Harnack inequality (Theorem 3.10) and some extensions are due to Serrin [SE 1]. This seems to be the first proof of a Harnack inequality by the maximum principle. Bers and Nirenberg [BN] derive a very similar result by altogether different and deeper methods.

The Liouville theorem (Corollary 3.12) is related to Bernstein's geometric theorem on surfaces of non-positive curvature (see [HO 4]) which implies that an entire solution $u$ of any elliptic equation $au_{xx} + 2bu_{xy} + cu_{yy} = 0$ such that $u = o(r)$ as $r \to \infty$ must be constant. Of particular interest is the fact that the equation need only be pointwise elliptic. In this generality Corollary 3.12 ceases to be valid, as counterexamples show. Bernstein's result is also based on the maximum principle but the argument is quite different and is more geometric.

\section*{Problems}

\begin{enumerate}
\item[3.1.] Let $L$ satisfy the conditions of Theorem 3.6 in a bounded domain $\Omega$ and suppose $Lu=0$ in $\Omega$.
\begin{enumerate}
\item[(a)] Let $\partial\Omega = S_1 \cup S_2$ ($S_1$ non-empty) and assume an interior sphere condition at each point of $S_2$. Suppose $u \in C^2(\Omega) \cap C^1(\Omega \cup S_2) \cap C^0(\overline{\Omega})$ satisfies the mixed boundary condition
\[
u = 0 \text{ on } S_1, \qquad \sum \beta_i D_i u = 0 \text{ on } S_2
\]
where the vector $\beta(x) = (\beta_1(x), \ldots, \beta_n(x))$ has a non-zero normal component (to the interior sphere) at each point $x \in S_2$. Then $u = 0$.
\item[(b)] Let $\partial\Omega$ satisfy an interior sphere condition, and assume that $u \in C^2(\Omega) \cap C^1(\overline{\Omega})$ satisfies the regular oblique derivative boundary condition
\[
\alpha(x)u + \sum \beta_i(x)D_i u = 0 \text{ on } \partial\Omega,
\]
where $\alpha(\beta \cdot \nu) > 0$, $\nu = \text{outward normal}$. Then $u = 0$.
\end{enumerate}
\item[3.2.] (a) If $L$ is elliptic, $Lu \ge 0$ ($\le 0$) and $c < 0$ in a domain $\Omega$, then $u$ cannot achieve an interior positive maximum (negative minimum). (No assumption is made concerning the coefficients $b^i$.)
\end{enumerate}

\source{gilbarg-trudinger:page-0060}

(b) If $L$ is elliptic with $c<0$ in a bounded domain $\Omega$, and $u\in C^2(\Omega)\cap C^0(\bar\Omega)$ satisfies $Lu=f$ in $\Omega$, then
\[
\sup_{\Omega}|u|\leq \sup_{\partial\Omega}|u|+\sup_{\Omega}|f/c|.
\]

\textbf{3.3.} Let $Lu=a u_{xx}+2b u_{xy}+c u_{yy}=0$ in an exterior domain $r>r_0$, $L$ being uniformly elliptic. Prove that if $u$ is bounded on one side then $u$ has a limit (possibly infinite) as $r\to\infty$; (cf. [GS]). [Apply the Harnack inequality in suitable annuli extending to infinity.] Use this result to prove the Liouville theorem, Corollary 3.12.

\textbf{3.4.} Let $u$ be a non-negative solution of
\[
Lu=a^{ij}D_{ij}u+b^iD_iu+cu=0,\qquad c\leq 0,\qquad i,j=1,2,
\]
where the coefficients of $L$ satisfy the inequalities
\[
\Lambda/\lambda\leq\mu,\qquad |b^i|/\lambda,\ |c|/\lambda\leq \nu\qquad (\mu,\nu=\text{const.}).
\]
Prove the Harnack inequality (3.21) with $K=K(\mu,\nu)$ and Corollary 3.11 with $\kappa=\kappa(\mu,\nu,\Omega,\Omega')$.

\textbf{3.5.} Assume the conditions on $L$ in Problem 3.4 are satisfied in the punctured disk $D_0:0<r<r_0$, and let $Lu=0$ in $D_0$. Prove that if $u$ is bounded on one side, then $u$ has a limit (possibly infinite) as $r\to 0$; (cf. [GS]).

\textbf{3.6.} Let $u\in C^2(\Omega)\cap C^0(\bar\Omega)$ be a solution of
\[
Lu=a^{ij}D_{ij}u+b^iD_iu+cu=f,\qquad c\leq 0
\]
in a bounded $C^1$ domain $\Omega$ of $\mathbb{R}^n$ satisfying an exterior sphere condition at $x_0\in\partial\Omega$, with $\bar B_R(y)\cap\bar\Omega=x_0$, and let $\lambda,\Lambda$ be positive constants such that
\[
a^{ij}(x)\xi_i\xi_j\geq \lambda|\xi|^2\qquad \forall x\in\Omega,\ \xi\in\mathbb{R}^n
\]
\[
|a^{ij}|,\ |b^i|,\ |c|\leq \Lambda.
\]
If $\varphi\in C^2(\bar\Omega)$ and $u=\varphi$ on $\partial\Omega$, show that $u$ satisfies a Lipschitz condition at $x_0$,
\[
|u(x)-u(x_0)|\leq K|x-x_0|,\ x\in\Omega,
\]
where $K=K(\lambda,\Lambda,R,\operatorname{diam}\Omega,\sup_\Omega|f|,\|\varphi\|_{2;\Omega})$. Hence conclude that $K$ provides a gradient bound for $u$ on $\partial\Omega$ when $u\in C^1(\bar\Omega)$ and $\partial\Omega$ is sufficiently smooth; (cf. [CH], p. 343). If the sign of $c$ is unrestricted, show that the same result holds provided $K$ depends also on $\sup_{\Omega}|u|$.

\source{gilbarg-trudinger:page-0061}

\textbf{Problems}

\textbf{3.7.} \quad (a) Let the operator $L$ in the preceding problem have H\"older continuous coefficients $a^{ij}$ at the origin: $|a^{ij}(x)-a^{ij}(0)|\leq K|x|^{\alpha}$, $\alpha>0$, in $|x|<r_0$ for some constant $K$. Suppose $Lu\geq 0$ $(c=0)$ in the punctured ball $0<r\leq r_0$, and assume
\[
u=
\begin{cases}
o(\log r), & n=2 \\
o(r^{2-n}), & n>2
\end{cases}
\qquad \text{as } r\to 0.
\]
Show that
\[
(3.31)\qquad \limsup_{x\to 0} u(x)\leq \sup_{|x|=r_0} u(x)
\]
and that equality holds only if $u$ is constant.

\quad (b) If $n>2$ show that the same conclusion holds as in part (a) if the coefficients $a^{ij}$ are continuous at $x=0$ and $u=o(r^{2-n+\delta})$ as $r\to 0$ for some $\delta>0$; (cf. [GS]).

\textbf{3.8.} Consider the equation
\[
(3.32)\qquad L_nu\equiv a^{ij}D_{ij}u=0,\qquad a^{ij}=\delta^{ij}+g(r)\frac{x_ix_j}{r^2},\qquad i,j=1,\ldots,n.
\]
Show that $L_nu=0$ has a radially symmetric solution $u=u(r)$, $r\neq 0$, satisfying the ordinary differential equation
\[
\frac{u''}{u'}=\frac{1-n}{r(1+g)}.
\]

\quad (a) If $n=2$ and $g(r)=-2/(2+\log r)$, show that equation (3.32) is uniformly elliptic in the disk $D:0\leq r\leq r_0=e^{-3}$ with continuous coefficients at the origin, and has bounded solutions $a+b/\log r$ in the punctured disk $D-\{0\}$ that do not satisfy (3.31).

\quad (b) If $n>2$ and $g(r)=-[1+(n-1)\log r]^{-1}$, show that (3.32) is uniformly elliptic in $0\leq r\leq r_0=e^{-1}$ and has continuous coefficients at the origin. Show that the corresponding solution $u=u(r)$ satisfies the condition $u=o(r^{2-n})$ as $r\to 0$ but does not satisfy (3.31).

\quad (c) If $n>2$, determine a function $g(r)$ such that (3.32) is uniformly elliptic and has a bounded solution $u=u(r)$ continuous at $r=0$ that does not satisfy (3.31).

\textbf{3.9.} \quad Let $w=z\,\exp\!\left[-(\log (1/|z|))^{1/2}\right]$. By considering the relation $w_{\bar z}=v(z)w_z$, where
\[
\frac{\partial}{\partial z}=\frac{1}{2}\left(\frac{\partial}{\partial x}-i\frac{\partial}{\partial y}\right),\qquad
\frac{\partial}{\partial \bar z}=\frac{1}{2}\left(\frac{\partial}{\partial x}+i\frac{\partial}{\partial y}\right),
\]
\source{gilbarg-trudinger:page-0062}
show that $u=\Re w=x\exp\bigl[-(\log(1/r))^{1/2}\bigr]$ satisfies a uniformly elliptic equation of divergence form
\[
(au_x+bu_y)_x+(bu_x+cu_y)_y=0,
\]
in which the coefficients $a\to 1$, $b\to 0$, $c\to 1$ at the origin and are regular in $0<r<1$. Observe that $u(0,0)=0$, $u(x,y)>0$ for $x>0$ and $u_x(0,0)=0$. Compare with Lemma 3.4.

\textbf{3.10.} Let $L$ be the operator of Problem 3.6, but without any condition on the sign of the coefficient $c$. Assume there is a function $v$ such that $v>0$ and $Lv\leq 0$ in $\Omega$. Then, if $Lu\geq 0$, show that the function $w=u/v$ cannot achieve a non-negative maximum in the interior of $\Omega$ unless it is constant.

\source{gilbarg-trudinger:page-0063}

\textbf{Chapter 4}

\begin{center}
\textbf{Poisson's Equation and the Newtonian Potential}
\end{center}

In Chapter 2 we introduced the fundamental solution $\Gamma$ of Laplace's equation given by
\begin{equation}
\Gamma(x-y)=\Gamma(|x-y|)=
\begin{cases}
\dfrac{1}{n(2-n)\omega_n}|x-y|^{2-n}, & n>2,\\[1.2ex]
\dfrac{1}{2\pi}\log |x-y|, & n=2.
\end{cases}
\tag{4.1}
\end{equation}

For an integrable function $f$ on a domain $\Omega$, the \textit{Newtonian potential} of $f$ is the function $w$ defined on $\R^n$ by
\begin{equation}
w(x)=\int_{\Omega}\Gamma(x-y)f(y)\,dy.
\tag{4.2}
\end{equation}

From Green's representation formula (2.16), we see that when $\partial\Omega$ is sufficiently smooth a $C^2(\OmegaBar)$ function may be expressed as the sum of a harmonic function and the Newtonian potential of its Laplacian. It is not surprising therefore that the study of Poisson's equation $\Delta u=f$ can largely be effected through the study of the Newtonian potential of $f$. This chapter is primarily devoted to the estimation of derivatives of the Newtonian potential. As well as leading to existence theorems for the classical Dirichlet problem for Poisson's equation, these estimates form the basis for the Schauder or potential theoretic approach to linear elliptic equations treated in Chapter 6.

\textbf{4.1. H\"older Continuity}

If the function $f$ in (4.2) belongs to $C_0^\infty(\Omega)$, we see by writing
\[
w(x)=\int_{\Omega}\Gamma(x-y)f(y)\,dy=\int_{\R^n}\Gamma(x-y)f(y)\,dy
\]
\[
=\int_{\R^n}\Gamma(z)f(x-z)\,dz
\]
